%%%% Modèle proposé par henry.costedoat@universite-paris-saclay.fr %%%%
%%%% Grand fan de LaTeX %%%%
%%%% Inspiré par un modèle proposé par frederic.mazaleyrat@ens-paris-saclay.fr 

%%%% 4 fevrier 2022  %%%%

\documentclass[french,12pt,a4paper]{report}
\usepackage[utf8]{inputenc}
\usepackage[T1]{fontenc}
\usepackage[french]{babel}
\usepackage{amsmath}
\usepackage{amsfonts}
\usepackage{fancyhdr}
\usepackage{amssymb}
\usepackage{color} % où xcolor selon l'installation
\definecolor{BleuFonce}{RGB}{0,74,117}
\usepackage{mdframed}
\usepackage{multirow} %% Pour mettre un texte sur plusieurs rangées
\usepackage{multicol} %% Pour mettre un texte sur plusieurs colonnes
\usepackage{tikz}
\usepackage{graphicx}
\usepackage[absolute]{textpos} 
\usepackage{colortbl}
\usepackage{array}
\usepackage{geometry}
\usepackage{hyperref}
\pagestyle{fancy}
\renewcommand\headrulewidth{1pt}
\usepackage{float}

\fancyfoot[L]{Polytech - - - - -} %% bas de page
\fancyhead[R]{Titre de rapport}  %% en-tete de page
\fancyhead[L]{Prénom Nom}

\begin{document}

\begin{titlepage}


%\thispagestyle{empty}

\newgeometry{left=6cm,bottom=2cm, top=1cm, right=1cm}

\tikz[remember picture,overlay] \node[opacity=1,inner sep=0pt] at (2.2mm,-165mm){\includegraphics{Fond1.png}}; % Fond changeable 

% fonte sans empattement pour la page de titre
\fontfamily{fvs}\fontseries{m}\selectfont

\color{white}

%%Texte vertical à gauche
\begin{picture}(0,0)
\put(-110,-743){\rotatebox{90}{\Huge{Nom du rapport / projet}}}
\end{picture}
 
%**************************************************************
%********************  LOGO  DE  POLYTECH  ********************
%****** CHANGER L'IMAGE POUR UN AUTRE POLYTECH QUE SACLAY *****
%* VOIR LES LOGOS DISPONIBLES, ET REMPLACER LE NOM CI-DESSOUS *
%**************************************************************

\vspace{-10mm} % à ajuster en fonction de la hauteur du logo
               % décalera le texte en dessous
\flushright \includegraphics[width=50mm]{Polytech_Saclay.png} % nom du logo du polytech




%*****************************************************
%******************** TITRE **************************
%*****************************************************
\flushright
\vspace{10mm} % à régler éventuellement
\color{BleuFonce}
\fontfamily{cmss}\fontseries{m}\fontsize{22}{26}\selectfont
 Titre du rapport / projet

\normalsize
\color{black}
%*****************************************************

%\fontfamily{fvs}\fontseries{m}\fontsize{8}{12}\selectfont

\vspace{1.5cm}
\normalsize

\textbf{Professeurs référents : M.Dupont et Mme.Martin}

\vspace{15mm}

Polytech - - - \\  %% nom école
\small Spécialité - - - \\  %% nom spécialité / enseignement

\vspace{15mm}

\textbf{Rapport de projet pour le 01/01/2022 à 23h59,\\ soutenu le 01/01/2022}\\
\bigskip
\Large {\color{BleuFonce} \textbf{Prénom Nom}}



\end{titlepage}

%%%%%%%%%%%%%%%%%%%%%%%%%%%%%%%%%%%%%%%%%%%%%%%%%%%%%%%%%%%%%%
\newpage

\setcounter{page}{1} 

\chapter*{Remerciements (chapitre non-numéroté)}

Texte 

\chapter{Nom chapitre}
\section{Nom sous-partie 1}

\par Texte sous partie 1

\section{Nom sous-partie 2}

\par Texte sous partie 2

\newpage
\chapter{Nom chapitre}
\section{Nom sous-partie 1}

\newpage
.
\newpage
.

\fontfamily{cmss}\fontseries{m}\selectfont

\small




%************************************
\vspace{\fill} % ALIGNER EN BAS DE PAGE
%************************************

% Modifier en fonction de l'école
\noindent
\color{BleuFonce} \footnotesize 
Polytech Paris-Saclay \\
Maison de l'Ingénieur\\
Bâtiment 620 Université Paris-Saclay\\
91190 Gif-sur-Yvette, France


\end{document}